\section{Introduction}
\label{sec:intro}

Knowledge graphs (KGs) represent facts as structured triples $(h, r, t)$ and underpin applications ranging from question answering and recommender systems to clinical decision support and drug repurposing~\cite{shen2022comprehensive,biswas2023knowledge}.
As KGs are increasingly deployed in safety-critical domains---pharmacovigilance, clinical decision support, and financial compliance---the gap between aggregate ranking metrics and per-prediction reliability has become consequential.
When these systems rely on KG completion to infer missing facts, they face a critical question: which predicted triples represent genuine knowledge, and which are merely statistical guesses?

\emph{Link prediction}---the task of inferring missing facts from the existing graph---has become a central problem in KG engineering and has spurred the development of a large family of embedding models~\cite{bordes2013translating, sun2019rotate, trouillon2016complex, dettmers2018convolutional, balazevic2019tucker, nguyen2018convkb, sun2022hake, le2023knowledge}.
The standard evaluation protocol for these models is straightforward and widely adopted: researchers compute Mean Reciprocal Rank (MRR) and Hits@$K$, two metrics that measure how highly the correct tail entity is ranked among all candidates. These point-estimate metrics have enabled reproducible benchmarking and supported rapid methodological progress over the past decade. Comprehensive surveys have catalogued the breadth of KG embedding research and increasingly call attention to the evaluation bottleneck~\cite{wang2017knowledge,rossi2021knowledge,li2023knowledge,liu2022understanding,biswas2023knowledge}.

Yet this evaluation paradigm carries a structural limitation: it condenses model behavior over thousands of test triples into a single scalar, discarding information about \emph{which} predictions are trustworthy and \emph{what fraction} of the test set the model handles reliably.
A model reporting MRR $= 0.56$ may appear capable, but this aggregate says nothing about whether it genuinely captures most relation types or performs near chance on a substantial subset of them.
Recent surveys have drawn attention to the lack of statistical rigor in KG completion evaluation~\cite{jarnac2025uncertainty,biswas2023knowledge}, identifying uncertainty quantification and reliability assessment as open problems.
The downstream stakes are material: applications such as drug repurposing and financial risk monitoring depend on KG completion, yet they cannot currently distinguish a confident correct prediction from a guess that happened to succeed.
Consider a pharmaceutical researcher using KG completion to prioritize candidate drug--target interactions for experimental validation. MRR = 0.56 tells them the model ranks well on average; it says nothing about whether the top-10 candidates for a specific relation type are 90\% reliable or 10\% reliable. This distinction directly affects experimental resource allocation---and it is invisible to current evaluation protocols.

We address this gap by demonstrating that KG completion evaluation can be profitably reformulated as a large-scale multiple hypothesis testing problem. Rather than introducing new FDR methodology, we adapt well-established statistical procedures---Benjamini-Hochberg FDR control~\cite{benjamini1995controlling}, Storey's $q$-value estimation~\cite{storey2002direct}, and Efron's empirical Bayes local FDR~\cite{efron2004large}---to the specific structure of KG link prediction.
This perspective aligns with knowledge-driven methodologies central to the KBS community~\cite{shen2022comprehensive,sun2022hake,le2023knowledge,zhang2023conflict,shang2024learnable,lu2024deep,purohit2025vanilla,fang2025knowledge}: the FDR framework provides a principled mechanism for distinguishing genuine knowledge discoveries from statistical artifacts, directly serving the needs of knowledge-based systems that consume KG completion output.
Our central technical observation is that link prediction can be recast as a large-scale multiple hypothesis testing problem: each test triple corresponds to a null hypothesis $H_0$ asserting that the triple is absent from the true KG, and the embedding model assigns each candidate a score that serves as a test statistic.
This reframing brings the statistical methods to bear on the problem, yielding principled, per-prediction reliability statements.
The framework is compatible with major embedding paradigms: it operates as a post-processing step on any trained KG embedding model and requires only the model's scoring function.

Together, these contributions establish a statistical evaluation paradigm that shifts the KG completion evaluation question from ``how well does the model rank on average?'' to ``for which specific predictions, relations, and models do we have statistical evidence of genuine knowledge discovery?''

We make the following contributions:

\begin{enumerate}
    \item We formalize KG completion evaluation as a multiple hypothesis testing problem, supplying a statistical foundation for assessing prediction reliability that augments the traditional ranking-based paradigm (Section~\ref{sec:method}).

    \item We develop an empirical $p$-value construction procedure that converts uncalibrated KG embedding scores into valid $p$-values via permutation-based negative sampling. The method imposes no distributional assumptions and allows FDR control procedures to be applied directly (Section~\ref{subsec:pvalue}).

    \item We introduce a three-level FDR-based evaluation protocol---global diagnostics, relation-stratified analysis, and cross-model comparison---that identifies dimensions of model quality invisible to scalar metrics (Section~\ref{subsec:protocol}).

    \item We evaluate the framework on WN18RR and FB15k-237 with TransE, RotatE, ComplEx, and ConvE. Our experiments yield three findings relevant to KG evaluation practice: (i) 26--85\% of test triples are statistically indistinguishable from noise ($\hat{\pi}_0 = 0.259$--$0.851$ depending on model and dataset), (ii) per-relation FDR rejection rates span $9.1\%$--$96.5\%$ under a single model, exposing extreme heterogeneity hidden by MRR, and (iii) MRR and FDR produce conflicting model rankings---on WN18RR, ConvE trails TransE on MRR ($0.469$ vs.\ $0.560$) yet achieves $2.8\times$ the BH rejection rate at $\alpha = 0.05$; the ranking divergence generalizes to FB15k-237 with different specific ordering (Section~\ref{sec:experiments}).
\end{enumerate}
